\section{Notations}\label{app-notation}
\begin{center}\small
\begin{tabular}{lp{12cm}}
\toprule
Symbole & Signification\\
\midrule
$S$, $r$, $R=rt$ & spot, taux zéro-coupon, taux cumulé\\
$b$, $B=bt$ & repo zéro-coupon, repo cumulé (portage)\\
$G=e^{R-B}$ & facteur de croissance du portage\\
$t_i$, $d_i$, $\tilde d_i=d_i/G(t_i)$ & dates de détachement, dividendes en numéraire, dividendes actualisés au portage\\
$D(T)$ & dividende cumulé actualisé au portage, \cref{eq-forward}\\
$T_k$, $k=1..K$ & échéances Totem (piliers)\\
$L_y$, $\pi_i$, $A_{ky}$ & niveau de dividendes de la période $y$, poids du profil, matrice d'allocation\\
$\Aset$, $\Bset$, $n_A$, $n_B$, $m$ & panels options et dividendes, leurs tailles, contributeurs communs aux deux\\
$\bA$, $\DA$, $\DB$ & moyennes des vues de repo et de dividendes du panel options ; consensus du panel dividendes\\
$y_k$ & écart de portage Totem moins FO, \cref{eq-y}\\
$X=[T,T\ln T]$, $\beta$ & modèle nul (écart de repo lisse) et ses coefficients\\
$\Sigma$, $\kappa$ & covariance du bruit de consensus, échelle de variance\\
$t$, $\tb$, $z_t$, $Q$ & gabarit, sa projection blanchie, statistique de filtre adapté, chi-deux global\\
$s_b$, $s_D$, $J$, $g$ & sensibilités FO au repo et aux dividendes, jacobienne du repo implicite, exposition non-forward\\
\bottomrule
\end{tabular}
\end{center}

\section{Démonstrations et remarques complémentaires}\label{app-proofs}

\begin{lemma}[Actualisations hétérogènes]
Si le contributeur $j$ actualise avec $DF_j$, le zéro de $\overline{C}-\overline{P}$ vaut $K^\ast=\sum_jDF_jF_j/\sum_jDF_j$,
une moyenne des forwards des contributeurs pondérée par l'actualisation. Son écart à la moyenne simple vaut
$\Cov(DF_j,F_j)/\overline{DF}$, du second ordre.
\end{lemma}
\begin{proof}
$\overline{C}-\overline{P}=\frac1n\sum_jDF_j(F_j-K)$ s'annule en $K^\ast$. Ensuite
$K^\ast-\overline F=\bigl(\overline{DF\,F}-\overline{DF}\,\overline F\bigr)/\overline{DF}$.
\end{proof}

\begin{lemma}[Rapport signal sur bruit exact dans un plan apparié]
Si un scénario ajoute un effet $\eta$ à l'écart de portage d'un sous-jacent sans modifier son bruit de consensus (même
graine), alors $y^{\text{scénario}}-y^{\text{sain}}=\eta$ et le rapport signal sur bruit $\|PL^{-1}\eta\|$ du
\cref{cor-power} est connu exactement, sans bruit. Pour un seul gabarit ($\eta=a\,t$), il vaut $|a|\,\|\tb\|$.
\end{lemma}
\begin{proof}
Le bruit s'élimine dans la différence ; pour $\eta=at$, $PL^{-1}\eta=a\tb$.
\end{proof}

\begin{lemma}[Calibration de Student du test d'erreur de marquage FO]
Soit $\hat\beta_0\sim\N(\beta_0,c\,\sigma^2)$ et $\hat\sigma^2$ une estimation indépendante avec
$\nu\hat\sigma^2/\sigma^2\sim\chi^2_\nu$. Alors $(\hat\beta_0-\beta_0)/\sqrt{c\hat\sigma^2}\sim t_\nu$.
\end{lemma}
\begin{proof}
C'est la définition de la loi de Student : une normale centrée réduite divisée par la racine d'un $\chi^2_\nu/\nu$
indépendant.
\end{proof}
Avec des dispersions estimées sur $n_A$ contributeurs, $\nu\approx n_A-1$ sur le court terme, là où se lit l'incertitude
du niveau de repo ; l'échelle de niveau $s$ corrige ensuite ce que cette approximation laisse échapper.

\begin{remark}[Terme du second ordre sur plusieurs piliers]
La \cref{prop-second} s'étend pilier par pilier : pour $V=\Phi(F_1,\dots,F_K)$ avec des échéances aux piliers, la dépendance
du total au premier ordre vis-à-vis du panel dividendes vaut
$-\sum_k\partial_k\Phi\,F_k\bigl[\tfrac12T_k^2\Delta b_k^2+T_k\Delta b_k\Delta D_k/(S-D_k)\bigr]$, plus les termes croisés créés
par l'interpolation du portage entre piliers, tous du second ordre.
\end{remark}

\section{Spécification du marché synthétique}\label{app-synthetic}
\begin{itemize}
\item \textbf{Dividendes.} Actions UE : dividende annuel détaché entre le 20 avril et le 30 mai ($35\,\%$ des titres versent
aussi un acompte de $40\,\%$ en novembre) ; actions US : trimestriel ; titres difficiles à emprunter : annuel ou
trimestriel ; indice UE : flux mensuel avec des poids $(1,2,4,14,24,18,6,4,3,4,3,2)$ par mois civil. Rendement et
croissance tirés par sous-jacent dans les intervalles du tableau des catégories.
\item \textbf{Repo.} $b(T)=b_\infty+(b_0-b_\infty)e^{-T/\tau}$, avec $b_0$, $b_\infty$ et $\tau$ tirés par sous-jacent.
\item \textbf{Vues dividendes des contributeurs.} $d_{ij}=d_i\bigl(1+\sigma(t_i)[0{,}7\,z_j+\sqrt{1-0{,}7^2}\,e_{ij}]\bigr)$, avec
$\sigma(t)$ affine en $t$ (propre à chaque catégorie), $z_j$ persistant et $e_{ij}$ idiosyncratique ; les dividendes annoncés
(avant $0{,}15$ an) ont un dixième de la dispersion.
\item \textbf{Vues repo des contributeurs.} $b_j(T)=b(T)+\sigma_b\bigl(u_j+0{,}6\,v_j\,\tau(T)\bigr)$, où $\tau(T)$ est une
inclinaison linéaire en $\ln T$ de $-0{,}5$ à $0{,}5$ ; $\sigma_b$ vaut $15$, $10$, $60$ et $8$\,pb pour les quatre catégories.
\item \textbf{Options.} Strikes de $80$ à $120\,\%$ du spot ; volatilité avec structure par terme et smile de $-0{,}10$ par
unité de log-moneyness, bruit de volatilité d'un point par contributeur, bruit de prix de $2\times10^{-5}$ du spot.
\item \textbf{Consensus.} Moyenne sur le panel après exclusion des cotations à plus de quatre écarts absolus médians
normalisés de la médiane.
\item \textbf{FO.} Dividendes avec un bruit de $1\,\%$, repo avec une erreur de niveau de $3$\,pb ; spot égal au spot de
référence Totem.
\end{itemize}

\section{Reproductibilité}\label{app-repro}
Tous les nombres, tableaux et figures de cette note sont produits par le dépôt :
\begin{center}\small
\begin{tabular}{l}
\code{pip install -e .[benchmark,report,dev]}\\
\code{pytest}\\
\code{python scripts/run\_benchmark.py} \quad (environ \vRuntime{} minutes)\\
\code{PYTHONPATH=. python scripts/build\_latex.py}
\end{tabular}
\end{center}
